% Generated draft: requires mathematical and source review.
\documentclass[11pt]{article}
\usepackage[margin=25mm]{geometry}
\usepackage{fontspec}
\setmainfont{Latin Modern Roman}
\usepackage{amsmath,amssymb,mathtools,unicode-math}
\setmathfont{Latin Modern Math}
\usepackage{enumitem,needspace,xurl,hyperref}
\setlength{\parindent}{0pt}
\setlength{\parskip}{5pt}
\setlength{\emergencystretch}{3em}
\newcommand{\N}{\mathbb{N}}
\newcommand{\Z}{\mathbb{Z}}
\newcommand{\Q}{\mathbb{Q}}
\newcommand{\R}{\mathbb{R}}
\newcommand{\C}{\mathbb{C}}
\newcommand{\F}{\mathbb{F}}
\newcommand{\A}{\mathbb{A}}
\newcommand{\PP}{\mathbb{P}}
\newcommand{\E}{\mathbb{E}}
\newcommand{\eps}{\varepsilon}
\newcommand{\dd}{\,\mathrm{d}}
\begin{document}
\section*{UnsolvedMath ID: 701343}
\textbf{Source problem code:} catalogue entry
\subsection{Definitions and mathematical statement}
For a group \(H\), write \(H^{(0)}=H\) and \(H^{(i+1)}=[H^{(i)},H^{(i)}]\). The group \(H\) is solvable if \(H^{(r)}=\{1\}\) for some integer \(r\geq 0\). For \(S\leq G\) and \(g\in G\), put \(S^g=g^{-1}Sg\).

Let \(G\) be a finite group satisfying
\[
N\trianglelefteq G\ \text{and}\ N\ \text{solvable}\quad\Longrightarrow\quad N=\{1\},
\]
and let \(S\leq G\) be solvable. The following questions are asked for every such pair \((G,S)\).
\begin{enumerate}
\item[(a)] (L. Babai, A. J. Goodman, L. Pyber.) Must there exist \(g_1,\ldots,g_7\in G\) such that
\[
\bigcap_{i=1}^{7}S^{g_i}=\{1\}?
\]
\item[(b)] Must there exist \(g_1,\ldots,g_5\in G\) such that
\[
\bigcap_{i=1}^{5}S^{g_i}=\{1\}?
\]
\end{enumerate}
Here the conjugating elements and the resulting conjugate subgroups are not required to be distinct. Thus each question equivalently asks for a trivial intersection of at most the stated number of conjugates. This is the conventional interpretation of the supplied wording; distinctness is not explicitly addressed in the source record.

The supplied editors' comments report a reduction of part (b) to almost simple groups, attributed to E. P. Vdovin (2012). A finite group \(H\) is almost simple if \(T\leq H\leq\operatorname{Aut}(T)\) for a nonabelian finite simple group \(T\); its socle is then \(T\). The comments report affirmative answers to part (b) for almost simple groups whose socles are isomorphic to an alternating group, a sporadic simple group, \(\operatorname{PSL}(n,q)\), \(\operatorname{PSU}(n,q)\), or \(\operatorname{PSp}(n,q)'\), where the prime denotes the derived subgroup. They also report an affirmative answer when \(S\) is a maximal subgroup of \(G\). The record does not specify parameter ranges or exceptional cases for the named classical families, nor the detailed formulation of the reduction. These comments and the registry's open status are supplied source assessments, not independently verified mathematical or literature conclusions.
\subsection{Short English statement}
Given a solvable subgroup \(S\) of a finite group \(G\) with no nontrivial solvable normal subgroup, must seven conjugates of \(S\) have trivial intersection? Must five suffice? The supplied editors' comments report a reduction and several affirmative cases for the five-conjugate question. This is a statement-only draft requiring human mathematical review.
\subsection{Sources}
\begin{itemize}
\item Local UnsolvedMath registry: supplied complete source record, permanent ID \(701343\), exact source problem code \texttt{catalogue entry}, title \emph{Kourovka Notebook Problem 17.41}. The record labels the problem open and attributes its proposal to E. P. Vdovin in 2010. Its supplied creation and update dates are 2026-10-16.
\item \emph{The Kourovka Notebook (Unsolved Problems in Group Theory)}, original problem 17.41, 17th issue (2010); supplied source edition: 21st issue (2026), \texttt{arXiv:1401.0300v48}. The collection record names E. I. Khukhro and V. D. Mazurov as editors. Part (a) is attributed in the statement to L. Babai, A. J. Goodman, and L. Pyber.
\item E. P. Vdovin, \emph{J. Algebra Appl.} \textbf{11}, no. 1 (2012), 1250015 (14 pages). Cited by the supplied editors for the reduction of part (b) to almost simple groups.
\item A. A. Baikalov, \emph{Algebra Logic} \textbf{56} (2017), 87--97. Cited by the supplied editors for the affirmative answer to part (b) for almost simple groups with alternating socle.
\item T. C. Burness, \emph{Israel J. Math.} \textbf{254} (2023), 313--340. Cited by the supplied editors for the affirmative answer to part (b) for almost simple groups with sporadic simple socle.
\item A. A. Baykalov, \emph{J. Group Theory} \textbf{28} (2025), 1003--1077. Cited by the supplied editors for the affirmative answer to part (b) for almost simple groups with socle isomorphic to \(\operatorname{PSL}(n,q)\).
\item A. A. Baykalov, \emph{Int. J. Algebra Comput.} \textbf{35}, no. 06 (2025), 823--908. Cited by the supplied editors for the affirmative answer to part (b) for almost simple groups with socle isomorphic to \(\operatorname{PSU}(n,q)\) or \(\operatorname{PSp}(n,q)'\).
\item T. C. Burness, \emph{Algebra Number Theory} \textbf{15}, no. 7 (2021), 1755--1807. Cited by the supplied editors for the affirmative answer to part (b) when \(S\) is maximal in \(G\).
\end{itemize}
\end{document}
